\scanpage{23}
\uline{Corollaire} $A$ anneau local \deleted{hensélien} noethérien
intègre, corps des fractions $K$, corps résiduel $k$.
Alors pour tt premier $r$, on a
\[
\mathrm{cd}_r(K)\geq\mathrm{cd}_r(k)+\dim A
\]
\deleted{[avec \uline{égalité} dans le cas $A$ hensélien \unsure{?}]}
\deleted{\inserted{\figure{\begin{tikzpicture}[baseline=(n.base)]\node[draw,ellipse,align=left,inner sep=3pt] (n) {et $A$ \unsure{Univ.} japonais\\\unsure{univ. caténaire}};\end{tikzpicture}}}}
\margin{\deleted{\unsure{à le}}}

sauf exception possible dans les cas suivants:

a) $r$ est la caractéristique de $K$\inserted{(égale caractéristique $>0$)}

b) $r=2$, \deleted{$A$ un anneau de \unsure{dim $1$}, $r=2$, Il existe}
\deleted{$A$ ``d'inégales caractéristiques'', \unsure{que} la caractéristique
résiduelle est $2$, \unsure{$r=$} un idéal premier $\mathfrak p\subset A$, \unsure{non nul}, avec $A/\mathfrak p$ de dim $1$.}
$\exists\ x\in\operatorname{Spec}(A)$ tel que $\kappa(x)$
soit ordonnable, et $\mathrm{cd}_2\kappa(x)(\sqrt{-1})$ soit fini
(corps réel \unsure{coulant|courant}).
\margin{\uline{Ex} $A=\mathbb R[x,y]/(x^2+y^2)$ \uline{\unsure{localisé}}}

\uline{NB.} Si $A$ est \uline{caténaire} et \uline{univ. japonais},
(ce que garantit le fait que les \unsure{lieux normaux}
des sous-schémas de $A$ sont ouverts), alors
\deleted{il suffit} on peut préciser\inserted{dans le cas b)}, \deleted{\unsure{à quel}} l'exception
possible par: $k$ lui-même ``corps réel \unsure{coulant|courant}''.
Enfin, lorsque $A$ est \uline{régulier},
cette exception elle-même est inoffensive, i.e.
on a l'inégalité annoncée, (pourvu qu'on ne
soit dans le cas a)].
\margin{cette restriction est-elle \unsure{inutile} ???}
\margin{22}

